\documentclass[11pt,a4paper]{article}
\usepackage{curso}
\begin{document}
\portada{Lección}{2.3}{La ecuación de Fokker--Planck}

\noindent\textbf{Unidad 2.}\quad\textbf{Objetivos:} evolución de la densidad; corriente de probabilidad; soluciones
estacionarias de equilibrio y fuera del equilibrio. \textbf{Prerrequisitos:} 2.2.

\section{De trayectorias a densidades}
Para $dx=\mu f(x)\,dt+\sqrt{2D}\,dW$, la densidad $p(x,t)$ del ensamble obedece
\begin{resultado}[ (Fokker--Planck)]
\begin{equation}\partial_tp=-\partial_xJ,\qquad J(x,t)=\mu f(x)\,p-D\,\partial_xp .\end{equation}
\end{resultado}
\begin{proof}[Derivación]
Para $g$ suave, Itô da $d\E g(x)=\E[g'\mu f+Dg'']\,dt$. Escribiendo $\E g=\int gp$ e integrando por partes:
$\int g\,\partial_tp=\int g\,[-\partial_x(\mu fp)+D\partial_x^2p]$ para toda $g$.
\end{proof}
Es una ecuación de continuidad: arrastre (velocidad $\mu f$) más difusión (de mayor a menor densidad). La animación
compara el histograma de 4000 trayectorias de Langevin \evS{} con la solución numérica de la ecuación (volúmenes
finitos, $\Delta x=0{,}01$, $\Delta t=10^{-4}$) en un pozo doble con $\kB T=0{,}35$.

\section{Estado estacionario y equilibrio}
Si $\partial_tp=0$, $J$ es constante. En 1D con $p\to0$ en los extremos (o fronteras reflectantes), $J=0$:
$\mu fp=D\partial_xp$. Con $f=-V'$ y $D=\mu\kB T$:
\begin{equation}p_{\text{eq}}(x)=\frac{e^{-V(x)/\kB T}}{Z}.\end{equation}
Esta es la demostración de que Langevin con la relación de Einstein muestrea Boltzmann (lección 2.2). En varias
dimensiones, el equilibrio exige $\mathbf J=0$ en todas partes (balance detallado).

\section{Estados estacionarios fuera del equilibrio (NESS)}
En un anillo con potencial periódico $U(\theta)=A\cos\theta$ y una fuerza constante $F$ (no derivable de un potencial
periódico), el estado estacionario tiene $J\neq0$: hay circulación neta. Con $A=1$, $F=0{,}8$, $D=0{,}5$ se obtiene
numéricamente $J\approx0{,}067$ \evS. La densidad no es Boltzmann y el sistema disipa $FJ\cdot2\pi$ por unidad de tiempo
(Unidad 3).

\section{Caso exacto: potencial cuadrático}
Para $V=\frac12\kappa x^2$ y $x(0)=x_0$:
\begin{equation}p(x,t)=\N\Big(x_0e^{-t/\tau},\ \tfrac{\kB T}{\kappa}(1-e^{-2t/\tau})\Big),\qquad\tau=\frac1{\mu\kappa}.\end{equation}
\begin{ejemplo}
$x_0=2$, $\kappa=1$, $\kB T=0{,}25$, $t=\tau$: $\langle x\rangle=2e^{-1}\approx0{,}74$, $\sigma^2=0{,}25(1-e^{-2})\approx0{,}22$.
\end{ejemplo}

\section{Espectro y mezcla}
El operador de Fokker--Planck $\mathcal L$ con $\partial_tp=\mathcal Lp$ tiene, en equilibrio, autovalores reales
$0=-\lambda_0>-\lambda_1\ge\dots$; la brecha $\lambda_1$ fija el tiempo de relajación (como $1-|\lambda_2|$ en la lección 0.5).
En el pozo doble, $\lambda_1$ es la suma de las tasas de Kramers de salto entre pozos.

\begin{aplicacion}
Corrección (Boltzmann como estacionaria), velocidad (brecha espectral), modelos de difusión (procesos directo e inverso
con score $\nabla\ln p$, Unidad 9) y producción de entropía $\dot\Sigma=\int J^2/(Dp)\,dx\ge0$ (Unidad 3).
\end{aplicacion}

\begin{error}
\begin{itemize}[nosep]
\item Estacionario $\neq$ equilibrio (NESS: $J\neq0$).
\item Condiciones de frontera (reflectantes, periódicas, absorbentes).
\item Ruido multiplicativo: la forma de Fokker--Planck depende de Itô/Stratonovich.
\item La ecuación describe la densidad, no una trayectoria.
\end{itemize}
\end{error}

\section{Resumen}
$\partial_tp=-\partial_xJ$, $J=\mu fp-D\partial_xp$; $J=0\Rightarrow$ Boltzmann; NESS con $J\neq0$; OU: solución gaussiana.

\begin{thebibliography}{9}\small
\bibitem[1]{gardiner} C. W. Gardiner, \emph{Handbook of Stochastic Methods}; H. Risken, \emph{The Fokker--Planck Equation}, 1996.
\bibitem[2]{peliti} L. Peliti y S. Pigolotti, \emph{Stochastic Thermodynamics: An Introduction}, 2021.
\bibitem[45]{song} Y. Song et al., Score-based generative modeling through SDEs, ICLR 2021.
\end{thebibliography}
\end{document}
